\chapter{Numerical Methods and Implementation Details}
\label{ch:implementation}

\section{Phase 1: constitutive bridge}

The network energy is evaluated in double precision on the CPU. PyTorch AD
computes \(\Pstress\) and \(\Aten\) from the scalar energy. NGSolve's
\texttt{IntegrationRuleSpace} stores deformation gradients and material
responses at matching quadrature points. Each Newton iteration:
\begin{enumerate}
  \item samples \(\F=\Iten+\nabla\bm u\) at material points;
  \item evaluates energy derivatives in PyTorch;
  \item writes \(\Pstress\) and the row-major \(9\times9\) tangent into
        integration-point grid functions;
  \item assembles the residual and stiffness; and
  \item solves the constrained Newton update.
\end{enumerate}
The quadrature rule attached to the integration-point space is reused in all
forms; otherwise the transferred material data would be evaluated at
different locations.

The deformation-gradient components are gathered from the integration-point
grid function into a point-major array of shape \((n_q,3,3)\). Before model
evaluation, the code checks finiteness and rejects any sampled \(J\leq10^{-8}\).
The constitutive interface then returns a first-Piola tensor and a material
tangent whose row and column indices both follow the \(3i+J\) flattening in
\cref{eq:flattening}. These values are written to nine- and 81-component
integration-rule grid functions, respectively. Keeping this mapping
explicit is essential: a transpose or a Voigt-style six-component shortcut
would assemble a different linearization.

At a current displacement \(u\), the assembled internal residual and tangent
are
\begin{equation}
  R_a(u)=\int_{\Omega_0} P_{iJ}\,N_{a,i,J}\,\dd V,
  \qquad
  K_{ab}(u)=\int_{\Omega_0}
  N_{a,i,J}A_{iJkL}N_{b,k,L}\,\dd V,
  \label{eq:bridge-residual-tangent}
\end{equation}
where \(N_{a,i,J}\) denotes the reference gradient associated with a vector
shape-function coefficient. The production solver constrains prescribed
boundary degrees of freedom and monitors the free residual. A backtracking
line search can use either residual reduction or the discrete energy;
orientation loss and non-finite constitutive values are rejected rather than
passed to the sparse solve. This separates invalid kinematics,
material-evaluation failure, and poor Newton globalization.

The Phase~1 training target is the compressible neo-Hookean energy in
\cref{eq:neo-hooke}. It can be evaluated for many deformation gradients
without generating a reference finite-element mesh or solving an RVE. The
network therefore isolates the PyTorch-to-NGSolve interface from
microstructure and homogenization error. The analytical material remains
available as a matched FE comparator after the bridge is operating.

\section{Periodic RVE and mixed matrix}

Opposing faces of a unit cube are identified geometrically in Netgen and
wrapped with NGSolve \texttt{Periodic} spaces, following the documented
periodic-mesh and periodic-space construction \cite{ngsolve_periodic}. A spherical particle is
represented by a cavity in the deformable matrix rather than by an
artificially stiff material. The fluctuation field uses \(P_2\) vector
elements and pressure uses \(P_1\) elements. This mixed approximation reduces
the locking risk for the nearly incompressible matrix; it does not make
locking independent of mesh, quadrature, or element choice. The nonlinear
residual is obtained from \cref{eq:mixed-energy}, and NGSolve supplies its
exact symbolic linearization.

The cell is the unit cube \([0,1]^3\), with a centered spherical cavity. Its
radius follows \(r=(3v_{\mathrm f}/4\pi)^{1/3}\); for one sphere in a
unit-volume cell this sets the particle-to-cell volume ratio. Netgen
identifies each pair of opposite faces before meshing, and NGSolve's
periodic spaces share the corresponding degrees of freedom. The macroscopic
deformation is prescribed as \(\Fbar\); the unknown field is the periodic
fluctuation \(w\), giving \(\F=\Fbar+\nabla w\). Periodicity therefore
applies to the fluctuation, not to the affine part of the total displacement.

The rigid particle is enforced through its cavity-boundary trace. The
prescribed trace is the exact finite rigid-body displacement
\(\bm t+(\Rten-\Iten)(\bm X-\bm X_c)\), with the particle center following
the macroscopic map. A free-particle solve determines the three rotational
coordinates by requiring zero resultant moment; translation is not solved
independently because a constant shift is the periodic gauge. Each rotation
update is multiplicative and uses the expansion
\(\exp([\delta q]_\times)=\Iten+[\delta q]_\times+
\frac12[\delta q]_\times^2+O(\|\delta q\|^3)\). If a direct initial load is
inadmissible, the solver can start at the identity state and proceed through
adaptive continuation, with alternative isochoric/volumetric paths and
previous-state or extrapolated predictors.

The mixed displacement-pressure equations are assembled from the energy
density using the same element integration order for the residual and its
symbolic linearization. In the standard dataset configuration, the matrix
shear scale is \(\mu=1\), bulk modulus \(\kappa=10^4\), displacement order is
two, and tetrahedral integration order is ten. These are dimensionless
constitutive scales for the numerical study; they do not represent an
experimentally calibrated PDMS parameter set.

\section{Rigid-particle equilibrium}

The particle trace is restricted to \cref{eq:rigid-motion}. Rotations are
updated multiplicatively using Rodrigues' formula. At each outer iteration,
the mixed RVE is equilibrated for the current rotation, the resultant moment
is assembled, and an exact condensed rotational Hessian supplies the Newton
step. Translation is not inverted as a rigid coordinate because it is the
periodic displacement gauge; the particle center is mapped by the
macroscopic deformation.

\section{Homogenized response}

The homogenized energy is the equilibrated mixed matrix potential per total
cell volume. Work-conjugate stress combines the matrix-volume contribution
with the full assembled particle reaction. The tangent augments the discrete
energy with nine global macroscopic coordinates and three incremental
rotation coordinates, then eliminates free displacement, pressure, and
rotation variables through \cref{eq:schur}. This calculation uses one sparse
factorization and multiple sensitivity solves; no dense inverse or numerical
stress perturbation is formed.

In practice, nine scalar global coordinates are added to the mixed FE space
to represent independent components of \(\Fbar\); three more represent
infinitesimal particle rotation when that degree of freedom is free. These
are global \texttt{NumberSpace} variables rather than nine additional
mesh-wide fields. After the equilibrium Hessian is assembled, the
macroscopic coordinates are held fixed and the free displacement, pressure,
and rotation block is factorized once. Each macroscopic unit perturbation is
then relaxed by solving the corresponding linear sensitivity problem. The
reaction recovered in the macroscopic coordinates supplies the tangent
column. The resulting tangent is checked for major symmetry and its action
is checked by Taylor remainders; finite perturbations are verification
tools, not the production definition of the tangent.

\section{Dataset generation and surrogate workflow}

Each macroscopic deformation state is assigned to exactly one split and is
evaluated at every requested volume fraction. This prevents deformation-state
leakage between training, validation, and test data. Individual RVE results
are checkpointed atomically before final NPZ assembly. The manifest records
configuration, field schema, software versions, shapes, and split counts.

The generated record contains
\begin{equation}
  \left(\vfrac,\Fbar,\overline\Psi,\Pbar,\Abar,\Rten,
  \text{diagnostics}\right).
\end{equation}
The energy, stress, and tangent supervision follows the Sobolev-training idea
of matching a function and its derivatives \cite{czarnecki2017}. The
context-conditioned model is
\texttt{HyperelasticNN(n\_ctx=1)}, an invariant-feature pICNN that is
isotropic in \(\Fbar\) and conditioned on volume fraction. The final
fine-tuning run used the Phase~6 normalization, energy and stress losses, and
a per-sample relative-Frobenius tangent loss with weights \(1\) and \(0.1\)
for stress and tangent, respectively. It used batches of 64, tangent
sub-batches of 8, learning rate \(10^{-4}\), seed 4103, and 3000 configured
epochs; the selected checkpoint was from epoch 2750. The frozen checkpoint
SHA-256 is
\texttt{32b0a41fdd75f7610a2147c053c013998f8cd6cb1032ea3e8c7d1b045c0fea70}.
Its grouped holdout energy and stress RMSEs are 0.519 and 2.466. This holdout
is not called a new independent evaluation because its Phase~7 deformation
states had already been used to evaluate the preceding checkpoint.

To improve simple-shear response, a separate 33-record direct-RVE dataset was
generated across \(\vfrac=0.02,0.05,0.10\), with 18 training, 6 validation,
and 9 test records grouped by deformation state. A full shear fine-tune
replayed the original Phase~6/7 data; a context-only adapter was also explored
as a development comparator. The selected full fine-tune was frozen before
the final holdout was generated. The final holdout uses a new microstructure
seed and four simple-shear amplitudes, with identity records retained only for
auditing. This final set was used only for evaluation, not model selection.

Standard deformation states are generated deterministically from a seeded
uniform perturbation of the identity and filtered to retain positive
determinant and a bounded condition number. Targeted compression and
expansion states can instead be stratified over specified determinant
intervals. A deformation group, rather than an individual volume-fraction
record, is assigned to train, validation, or test; all volume-fraction
variants of a state remain together. This avoids leakage of identical
loading states across splits. A solved sample is accepted only if the mixed,
rotational, and condensed residuals are below \(10^{-8}\), the sampled
microscopic Jacobian exceeds its configured floor, particle clearance is
positive, and all stored arrays are finite.

The schema-2 record stores the loading gradient and determinant, volume
fraction, homogenized energy, \(3\times3\) stress, row-major \(9\times9\)
tangent, optimized particle rotation, solver iteration counts, equilibrium
residuals, tangent-symmetry error, and minimum/maximum sampled microscopic
Jacobian. A JSON manifest accompanies the compressed array file and records
the configuration, field shapes and dtypes, software version, split counts,
and hashes. Per-sample temporary files are atomically renamed after a
successful solve, allowing a restarted generation job to keep completed
samples without treating a partially written result as valid.

Training is performed in float64. Feature, context, energy, stress, and
tangent scales are fitted from the training split only, or inherited from
the parent checkpoint during fine-tuning. The loss combines normalized
energy and stress errors; selected states also contribute a per-sample
relative-Frobenius tangent term. Tangent examples are sub-batched because
second derivatives are more memory intensive than energy evaluation.
Validation is grouped by deformation state, and the selected epoch is
determined from validation loss rather than the test set. Checkpoints retain
the model state, normalization, training configuration, and provenance
needed to recreate inference.

\section{Macroscopic finite-element solution}

The reported component-scale problem has a displacement field only. NGSolve
applies the macroscopic boundary conditions and solves for displacement;
PyTorch evaluates the homogenized energy, stress, and tangent at its
integration points. The pressure field belongs to the microscale mixed RVE
problem and is not an additional macroscopic unknown. The workflow is offline
homogenization followed by a surrogate macro solve, not concurrent FE\(^2\).

The trained surrogate is evaluated at macroscopic integration points and
returns energy, first-Piola stress, and the exact AD tangent. With first-order
displacements, fourth-order integration, an energy merit function, a secant
load predictor, and 100 load steps, the Phase~6/7 checkpoint reached the
target 5\% stretch at \(h=0.25,0.20,0.16\). Their final free-residual norms were
\(1.63\times10^{-10}\), \(3.79\times10^{-12}\), and
\(1.04\times10^{-12}\), respectively. Second-order runs did not reach the
target load. Consequently, load-path completion is demonstrated for the
first-order setting, while high-order stability, response convergence, and
component-level accuracy remain open.

At component scale, NGSolve imposes the external displacement boundary
conditions and assembles the macroscopic residual and stiffness. At each
quadrature point, the material bridge sends the current deformation gradient
and fixed volume-fraction context to the PyTorch pICNN, which supplies
energy, stress, and tangent. Newton updates are globalized with the discrete
energy merit function; a secant predictor initializes each prescribed load
increment from the two preceding converged states. The tested paths use
100 load steps to reach five-percent stretch. Pressure remains internal to
the RVE calculations and is not introduced as a second macroscopic field.
